% Part: second-order-logic
% Chapter: sol-and-sets
% Section: cardinalities

\documentclass[../../../include/open-logic-section]{subfiles}

\begin{document}

\olfileid[hi]{sol}{set}{crd}

\olsection{समुच्चयों के आकार}

\begin{explain}
जिस प्रकार हम व्यक्त कर सकते हैं कि प्रांत परिमित अथवा अनंत,
!!{enumerable} अथवा !!{nonenumerable} है, उसी प्रकार हम~$\Domain{M}$ के
किसी उपसमुच्चय के परिमित अथवा अनंत, !!{enumerable} अथवा !!{nonenumerable}
होने के गुण को परिभाषित कर सकते हैं।
\end{explain}

\begin{prop}
सूत्र $\fn{Inf}(X) \ident$
\begin{multline*}
\lexists[u][(\lforall[x][\lforall[y][(\eq[u(x)][u(y)] \lif 
      \eq[x][y])]] \land {} \\\lexists[y][(X(y) \land \lforall[x][(X(x)
      \lif \eq/[y][u(x)])]])]
\end{multline*}
किसी चर-नियतन~$s$ के सापेक्ष तभी और केवल तभी संतुष्ट होता है जब $s(X)$
अनंत हो।
\end{prop}

\begin{prop}
सूत्र $\fn{Count}(X) \ident $
\begin{multline*}
\lexists[z][\lexists[u][(X(z) \land 
    \lforall[x][(X(x) \lif X(u(x)))] \land {} \\ \lforall[Y][((Y(z) \land
      \lforall[x][(Y(x) \lif Y(u(x)))]) \lif X = Y])]])
\end{multline*}
किसी चर-नियतन~$s$ के सापेक्ष तभी और केवल तभी संतुष्ट होता है जब $s(X)$
!!{enumerable} हो।
\end{prop}

कैंटर के प्रमेय से हम जानते हैं कि !!{nonenumerable} समुच्चय होते हैं, और
वास्तव में अनंत आकारों के अनंत रूप से अनेक भिन्न स्तर होते हैं।
समुच्चय-सिद्धांत समुच्चयों के आकारों का पूरा अंकगणित विकसित करता है और
समुच्चयों को अनंत कार्डिनल संख्याएँ देता है। प्राकृतिक संख्याएँ परिमित
समुच्चयों के आकार मापने वाली कार्डिनल संख्याओं का काम करती हैं।
!!{denumerable} समुच्चयों की कार्डिनैलिटी पहली अनंत कार्डिनैलिटी है, जिसे
~$\aleph_0$ (``अलेफ़-नॉट'' अथवा ``अलेफ़-शून्य'') कहते हैं। अगला अनंत आकार
~$\aleph_1$ है। यह वह सबसे छोटा आकार है जो किसी समुच्चय का परिगणनीय
(अर्थात् आकार~$\aleph_0$ का) हुए बिना हो सकता है। ``$X$ का आकार
$\aleph_0$ है'' को हम $\fn{Aleph}_0(X) \liff \fn{Inf}(X) \land
\fn{Count}(X)$ से परिभाषित कर सकते हैं। $X$ का आकार $\aleph_1$ तभी और केवल
तभी है जब उसके सभी उपसमुच्चय परिमित हों या आकार~$\aleph_0$ के हों, किंतु
वह स्वयं आकार~$\aleph_0$ का न हो। अतः इसे सूत्र $\fn{Aleph_1}(X) \ident
\lforall[Y][(Y \subseteq X \lif (\lnot\fn{Inf}(Y) \lor
\fn{Aleph}_0(Y)))] \land \lnot \fn{Aleph}_0(X)$ द्वारा व्यक्त किया जा
सकता है। आकार $\aleph_2$ का होना भी इसी प्रकार परिभाषित होता है, इत्यादि।

एक आकार विशेष रुचि का है, जिसे सातत्य की कार्डिनैलिटी कहते हैं। यह
$\Pow{\Nat}$ का आकार है, अथवा तुल्य रूप से~$\Real$ का आकार। किसी समुच्चय का
आकार सातत्य के बराबर है, यह बात भी द्वितीय-कोटि तर्कशास्त्र में व्यक्त की
जा सकती है, किंतु इसके लिए कुछ अधिक काम करना पड़ता है।

\end{document}
